%! Polynomial Functions and Complex Zeros — Unit 2, Lesson 2.4 — Lesson Plan
\documentclass[10pt]{article}
\usepackage{precalculus-article}
\usepackage{precalculus-boxes}
\newcommand{\TallMath}[1]{$\displaystyle #1\rule[-1.4em]{0pt}{3.2em}$}
\newcommand{\CourseName}{Precalculus}
\newcommand{\MeetingLength}{60 minutes}
\newcommand{\UnitNumberName}{Unit 2: Polynomial Functions \quad}
\newcommand{\LessonNumberName}{Lesson 2.4: Polynomial Functions and Complex Zeros}

\begin{document}
\begin{center}
    {\Large\bfseries \CourseName \\
    {\small\bfseries \UnitNumberName \LessonNumberName}}
\end{center}

\begin{tcolorbox}[colback=lilac, colframe=plum, boxrule=0.9pt, arc=2mm,
    left=2mm, right=2mm, top=1.5mm, bottom=1.5mm]
    \textbf{Primary Objective:} Students will use the Fundamental Theorem of
    Algebra to count the complex zeros of a polynomial of degree $n$, reconcile
    that count with the $x$-intercepts visible on a graph, use conjugate pairs
    to name a missing non-real zero and to build a polynomial from a zero list,
    and classify a polynomial function as even, odd, or neither.
    \hfill\textit{\small (CED Topic 1.5 --- Skills 1.B, 2.B)}
\end{tcolorbox}

\renewcommand{\arraystretch}{1.6}
\begin{skillbox}[Priority Ideas \& Skills]{goldbox}
\begin{tabularx}{\linewidth}{@{}X X@{}}
\textbf{AP Skills} &
\textbf{Key Understandings} \\
\midrule
\textbf{1.B} --- Express functions, equations, or expressions in analytically
equivalent forms useful in a given context. &
\begin{itemize}[leftmargin=1.2em,itemsep=2pt,topsep=0pt]
  \item If $a$ is a complex number and $p(a) = 0$, then $a$ is a zero of $p$
        and a root of $p(x) = 0$. Zeros need not be real. (1.5.A.1)
  \item A polynomial function of degree $n$ has \textbf{exactly $n$ complex
        zeros}, counting multiplicities --- the Fundamental Theorem of
        Algebra. (1.5.A.2)
\end{itemize} \\
\textbf{2.B} --- Construct equivalent graphical, numerical, analytical, and
verbal representations of functions. &
\begin{itemize}[leftmargin=1.2em,itemsep=2pt,topsep=0pt]
  \item Real zeros are the $x$-intercepts you can see; non-real zeros are
        nowhere on the graph. Degree $-$ (real zeros with multiplicity)
        $=$ number of non-real zeros. (1.5.A.2)
  \item If $a + bi$ is a zero of a polynomial with real coefficients, so is
        $a - bi$. Non-real zeros arrive two at a time, so their count is even
        and an odd-degree polynomial always has a real zero. (1.5.A.4)
\end{itemize} \\
\textbf{LO 1.5.B} --- Determine whether a polynomial function is even or odd. &
\begin{itemize}[leftmargin=1.2em,itemsep=2pt,topsep=0pt]
  \item Even: $f(-x) = f(x)$, symmetric over the line $x = 0$. Odd:
        $f(-x) = -f(x)$, symmetric about the point $(0,0)$. (1.5.B.1--2)
  \item Shortcut for polynomials: all-even exponents $\Rightarrow$ even;
        all-odd exponents $\Rightarrow$ odd; a mix $\Rightarrow$ neither. A
        constant term counts as $x^0$, an even power.
\end{itemize} \\
\end{tabularx}
\end{skillbox}

\begin{skillbox}[{Vocabulary, Concepts, \& Theorems}]{greenbox}
\renewcommand{\arraystretch}{1.4}
\begin{tabularx}{\linewidth}{@{} >{\leavevmode\bfseries\color{plum}}p{0.27\linewidth} X @{}}
Imaginary unit    & The number $i$, defined by $i^2 = -1$. \\
Complex number    & Any number $a + bi$ with $a$ and $b$ real; $a$ is the real part and $b$ the imaginary part. Every real number is complex (with $b = 0$). \\
Non-real zero     & A zero $a + bi$ with $b \ne 0$. It is not an $x$-intercept and does not appear on the graph at all. \\
Complex conjugate & The partner $a - bi$ of $a + bi$. Their product $(a+bi)(a-bi) = a^2 + b^2$ is real. \\
Fundamental Theorem of Algebra & A polynomial function of degree $n$ has exactly $n$ complex zeros, counting multiplicities. \\
Even function     & $f(-x) = f(x)$ for every $x$ in the domain; the graph is symmetric over the line $x = 0$. \\
Odd function      & $f(-x) = -f(x)$ for every $x$ in the domain; the graph is symmetric about the point $(0,0)$. \\
\end{tabularx}
\end{skillbox}

\begin{skillbox}[Activate Prior Knowledge \& Spiral Review]{lilac}
    Please see attached warm-up (spiral review).
    Skills reviewed: reading real zeros and their multiplicities off a factored
    polynomial and totaling the degree from the exponents on its factors
    (Lesson 2.3), and one item of Algebra 2 complex arithmetic --- $i^2 = -1$
    and the product of a conjugate pair. Item 2's factor $x^2 + 1$ is the one
    students cannot break apart over the real numbers; that dead end is the
    door into today's lesson.
\end{skillbox}

\begin{skillbox}[Hook]{lilac}
    \textbf{The missing zeros.} Project the pre-drawn graph of
    $p(x) = x^4 - 3x^2 - 4$ from the notes. It is a degree-4 polynomial, and
    its graph crosses the $x$-axis exactly twice, at $x = -2$ and $x = 2$.
    \begin{enumerate}[itemsep=2pt]
        \item Ask: Lesson 2.3 tied zeros to factors, and a degree-4 polynomial
              has four factors' worth of zeros. We can see two. Are the other
              two missing, or are we not looking in the right place?
        \item Factor on the board: $p(x) = (x-2)(x+2)(x^2+1)$. The first two
              factors give the intercepts. Ask what $x^2 + 1 = 0$ gives.
    \end{enumerate}
    \textit{Discussion prompt:} $x^2 = -1$ has no \textit{real} solution ---
    but it has two solutions the moment we allow $i$. Today the count comes out
    exactly right, every time, provided we are willing to look off the graph.
\end{skillbox}

\begin{skillbox}[Lesson]{lilac}
\begin{multicols}{2}

\textbf{Part 1: A Complex-Number Refresher} (10 min)

\begin{itemize}
  \item Recall from Algebra 2: $i^2 = -1$. A \textbf{complex number} is
        $a + bi$ with $a$ and $b$ real --- real part $a$, imaginary part $b$.
  \item Every real number is already complex ($b = 0$). ``Non-real'' means
        $b \ne 0$; that is the new territory.
  \item Evaluate at a complex input exactly as always: for $p(x) = x^2 + 1$,
        $p(i) = i^2 + 1 = -1 + 1 = 0$. So $i$ \textit{is} a zero, by the same
        definition used in 2.3.
  \item Conjugates: $(3+2i)(3-2i) = 9 + 4 = 13$. The $i$ terms cancel and the
        product is real --- hold on to that; it explains Part 3.
\end{itemize}

\textbf{Part 2: Exactly $n$ Zeros} (13 min)

\begin{itemize}
  \item State it plainly: a polynomial function of degree $n$ has
        \textbf{exactly $n$ complex zeros, counting multiplicities}. Not ``at
        most'' --- exactly.
  \item Return to the hook: degree 4, zeros $2, -2, i, -i$. Four of them. The
        graph showed two because only real zeros can be $x$-intercepts.
  \item Fill the zero-census table together. The subtraction to rehearse:
        degree $-$ (real zeros counted with multiplicity) $=$ non-real zeros.
\end{itemize}

\columnbreak

\textbf{Part 3: Non-Real Zeros Come in Pairs} (12 min)

\begin{itemize}
  \item If $a + bi$ is a zero of a polynomial with real coefficients, then
        $a - bi$ is a zero too. Point back to Part 1: the conjugate pair is
        what multiplies out to a \textit{real} quadratic factor.
  \item Two consequences worth writing down: the number of non-real zeros is
        always \textbf{even}, and an \textbf{odd-degree} polynomial must have
        at least one real zero --- so its graph must cross the $x$-axis.
  \item Build from a zero list: zeros $2$ and $1 + 3i$ force $1 - 3i$, giving
        $(x-2)(x^2 - 2x + 10) = x^3 - 4x^2 + 14x - 20$. Point out that the
        coefficients came out real.
\end{itemize}

\textbf{Part 4: Even and Odd Polynomials} (12 min)

\begin{itemize}
  \item Definitions first: $f(-x) = f(x)$ is \textbf{even} (mirror over the
        line $x = 0$); $f(-x) = -f(x)$ is \textbf{odd} (half-turn about the
        origin).
  \item Verify by substitution, not by eye: for $g(x) = x^3 - 4x$,
        $g(-x) = -x^3 + 4x = -g(x)$.
  \item Then the polynomial shortcut: all-even exponents $\Rightarrow$ even,
        all-odd $\Rightarrow$ odd, a mix $\Rightarrow$ neither. Flag the trap:
        a constant is $x^0$, an even power, so $x^3 + 5$ is neither.
  \item Close the loop --- the hook polynomial $x^4 - 3x^2 - 4$ is even, and
        the pre-drawn graph shows the mirror.
\end{itemize}
\end{multicols}
\end{skillbox}

\begin{skillbox}[Active Monitoring]{redbox}
While students work on guided notes and practice, circulate and check:
\begin{itemize}
  \item \textbf{Common error:} reading ``4 zeros'' off a graph with 2
        intercepts. Cold-call: ``How many of those four does the graph get to
        show you? Where are the other two hiding?''
  \item \textbf{Common error:} dropping multiplicity in the census ---
        counting $(x+1)^3$ as one real zero instead of three. Ask: ``Counting
        \textit{multiplicities}. How many times does that factor appear?''
  \item \textbf{Common error:} giving the conjugate of $3 - 5i$ as $-3 + 5i$.
        Only the \textit{imaginary} part flips sign; the real part stays put.
  \item \textbf{Common error:} calling $x^3 + 5$ odd because of the $x^3$.
        Push: ``What power of $x$ is the 5 sitting on? Is that even or odd?''
  \item \textbf{Language:} require ``non-real zero,'' not ``imaginary zero''
        or ``fake zero.'' The zeros are perfectly real objects; they are just
        not real \textit{numbers}.
\end{itemize}
\end{skillbox}

\begin{skillbox}[Group Work \& Differentiation]{redbox}
Students work on the \textbf{Group Activity: The Missing Zeros Bureau} in
leveled tiers. Every tier reads the same pre-drawn cubic graph --- no tier
draws anything.

\begin{itemize}
  \item \textbf{Tier R (Remediate):} Total the zeros from the degree alone,
        subtract the real ones to get the non-real count, and name conjugate
        partners for four given non-real numbers.
  \item \textbf{Tier A (Approaching Proficiency):} Build a degree-3
        polynomial from the zeros $-1$ and $2i$ (supplying $-2i$), then
        classify four functions as even, odd, or neither and verify one by
        substitution.
  \item \textbf{Tier E (Extension):} Reconcile the pre-drawn graph of
        $p(x) = x^3 - 3x^2 + 4x - 2$ with the algebraic count --- one visible
        intercept, degree 3, so two non-real zeros --- confirm $x = 1$ by
        evaluation, find $1 \pm i$ from the remaining quadratic, and justify
        in writing why every odd-degree polynomial with real coefficients must
        cross the $x$-axis (AP Skill 2.B).
\end{itemize}
\end{skillbox}

\begin{skillbox}[Individual Work \& Assessment]{redbox}
\textbf{Exit Ticket} (last 5--7 minutes, individual, collected):
\begin{enumerate}
  \item For $p(x) = (x-4)(x^2+25)$: give the degree, list all complex zeros
        (worked space provided), and state how many $x$-intercepts the graph
        has.
  \item A degree-4 polynomial with real coefficients has $2 - i$ as a zero.
        Name another zero, then say how many real zeros it has, counting
        multiplicity.
  \item Classify three polynomials as even, odd, or neither.
\end{enumerate}
\textbf{Assessment note:} Item 1(c) is the diagnostic --- a student who writes
``3 intercepts'' because they found three zeros has not separated the
algebraic count from the graph, and that separation is the whole lesson. Item
3(c), $x^3 + x^2$, is the mixed-exponent check; a miss there is a definition
miss, not an arithmetic slip.
\end{skillbox}

\begin{skillbox}[Reinforcement \& Extension]{goldbox}
\textbf{Homework:} 5 problems covering degree and complete zero lists from
factored forms, a zero-census table completed by subtraction, building a
degree-3 polynomial from the zeros $5$ and $1 - 2i$, an even/odd/neither
classification set with one substitution verification, and one AP-style
multiple-choice item on counting non-real zeros.

\textbf{Extension (optional):} Multiply out
$\big(x - (a+bi)\big)\big(x - (a-bi)\big)$ to get $x^2 - 2ax + (a^2+b^2)$ ---
a quadratic with \textit{real} coefficients. That is why conjugate pairs are
forced, and why every real-coefficient polynomial factors into real linear and
real quadratic pieces.

\textbf{Preview:} Next lesson (2.5, Polynomial Functions and End Behavior)
leaves the zeros behind and looks at the two far ends of the graph: what the
degree and the leading coefficient together predict as $x$ runs off toward
$\infty$ and toward $-\infty$.
\end{skillbox}

\begin{teachernote}[Warm-Up]
Five minutes, silent start, no calculators. Items 1 and 2 are pure 2.3 recall
and should be quick; do not reteach multiplicity here, just confirm it. Item
2's $x^2 + 1$ factor is deliberately un-factorable over the reals --- if a
student announces ``that one has no zeros,'' write the claim on the board and
save it for the hook rather than correcting it now. Item 3 is the Algebra 2
hinge: a class that cannot produce $i^2 = -1$ on demand needs Part 1 stretched
to a full 15 minutes.
\end{teachernote}

\begin{teachernote}[Guided Notes]
Part 2 is the centerpiece and deserves the most air time --- ``exactly $n$,''
with the word \textit{exactly} circled. The census table is the backbone: work
row 1 aloud, let pairs do rows 2--4, then share. Row 1 is the hook polynomial,
so the graph on the previous page is the answer key for it. Part 3's
conjugate-pair result is stated, not proved; the proof is the homework
extension. Keep Part 1 tight --- it is a refresher, and the only two moves the
lesson needs from it are evaluating $p$ at a complex input and recognizing a
conjugate pair. If time runs short, compress Part 1, never Part 4: even/odd is
its own learning objective and appears on the exit ticket.
\end{teachernote}

\begin{teachernote}[Group Activity]
All three tiers reference the same cubic graph, so groups can be regrouped
mid-activity without new setup. Tier R is deliberately arithmetic-light: the
entire task is ``degree minus real equals non-real,'' repeated until it is
automatic. Tier A's build-a-polynomial item is the one to check for real
coefficients at the end --- if a student's expanded answer still contains an
$i$, they multiplied the conjugate factors wrong, and that is worth the whole
group's attention. Tier E's question 3 is the writing task; accept any version
of ``non-real zeros come in pairs, so an odd total cannot be all non-real.''
\end{teachernote}

\begin{teachernote}[Exit Ticket]
Collected, completion credit only. Sort by item 1(c) first: ``1 intercept'' is
the target, ``3 intercepts'' is the reteach pile. On item 2(b), accept ``2''
with or without reasoning; the reasoning surfaces in tomorrow's warm-up either
way. Item 3 should take under a minute for a student who internalized the
exponent shortcut --- if it took much longer, Part 4 needs a five-minute
revisit before 2.5.
\end{teachernote}

\begin{teachernote}[Homework]
Problem 2's census table is the one to grade first; it isolates the counting
idea from every other skill in the lesson. Problem 5 is the AP-style item ---
distractor D (``3'') is chosen by students who subtract without remembering
that non-real zeros come in pairs, and it is worth two minutes of autopsy
tomorrow because it is exactly the misconception 2.7 will punish. The
extension is the proof behind Part 3; if the class produced the ``why do they
have to come in pairs?'' question during the lesson, assign it explicitly and
open tomorrow with a student presenting it.
\end{teachernote}
\end{document}
